% Separate addition: does not replace 07_elipse.tex.
% Provenance: integral 2249, 52.4.1--2, 52.6.2--3, 53.28, and 90.4.6.
% Recovery of results and geometric clarification; local certificate
% verificar_elipse_II.py. Uses no physical radius as a generating input.

\subsection{Semiaxes, foci, and interior radii}
\label{sec:ii-elipse-geometria}

The figure is obtained from the angular and action outputs of the preceding
sections. APP preserves the sheets and their residues and quotients; TRIT fixes
orientation; TPK transports boundary, carry, and memory through to the
representations of the same enriched state and the joint discrete
structure of the continuum. The pair \((A,C^*)\) and the positive section
\(\hbar\) are received here after that construction. The geometry makes explicit
what information their readers preserve, without using an observed length
to select the source state.

Let \(\xi=C^*/A\in(-1,1)\), with \(A\ne0\) and both angles in the
same unit. Write \(\eta=\operatorname{artanh}\xi\) and
\[
 a_S=\sqrt{2\hbar}\,e^{\eta/2},\qquad
 b_S=\sqrt{2\hbar}\,e^{-\eta/2},\qquad
 \gamma(\theta)=(a_S\cos\theta,b_S\sin\theta).
\]
The semiaxes are labeled by their coordinates: if \(\eta<0\),
the vertical axis is the larger. Both belong to the square-root-of-action
line of the balanced chart, not to the line of spatial lengths.

\begin{proposition}[Focal geometry and radial interior]
\label{prop:ii-elipse-focos}
Let \(a=\max(a_S,b_S)\), \(b=\min(a_S,b_S)\). The focal half-distance
and eccentricity are
\begin{equation}
 f_S=\sqrt{a^2-b^2}
 =\sqrt{4\hbar\sinh|\eta|},\qquad
 e_f=\frac{f_S}{a}=\sqrt{1-e^{-2|\eta|}}.
 \label{eq:ii-elipse-focal}
\end{equation}
For \(\eta>0\), the foci are \((\pm f_S,0)\); for
\(\eta<0\), they are \((0,\pm f_S)\). At \(\eta=0\), they coincide with
the center. The boundary radius in polar direction \(\phi\) is
\begin{equation}
 \rho_{\partial}(\phi)
 =\left(\frac{\cos^2\phi}{a_S^2}
        +\frac{\sin^2\phi}{b_S^2}\right)^{-1/2}.
 \label{eq:ii-elipse-radio-polar}
\end{equation}
The interior segment in that direction consists of
\(r(\cos\phi,\sin\phi)\), with \(0\le r<\rho_{\partial}(\phi)\).
In particular, \(b\le\rho_{\partial}(\phi)\le a\).
\end{proposition}

\begin{proof}
The squares of the major and minor semiaxes are
\(2\hbar e^{|\eta|}\) and \(2\hbar e^{-|\eta|}\). Their difference
gives \eqref{eq:ii-elipse-focal}. To verify the focus, suppose
\(\eta>0\). At \(X=\gamma(\theta)\), the distances to
\(F_+=(f_S,0)\) and \(F_-=(-f_S,0)\) satisfy
\[
 |X-F_+|^2=(a_S-f_S\cos\theta)^2,\qquad
 |X-F_-|^2=(a_S+f_S\cos\theta)^2.
\]
Since \(a_S>f_S\), their positive roots sum to \(2a_S\).
For \(\eta<0\), exchanging coordinates gives the same proof.
In the circular case, \(f_S=0\). Finally, substituting
\((Q,P)=r(\cos\phi,\sin\phi)\) into
\(Q^2/a_S^2+P^2/b_S^2=1\) gives
\eqref{eq:ii-elipse-radio-polar}. The coefficient of \(r^2\) lies
between \(a^{-2}\) and \(b^{-2}\), proving the bounds and the
characterization of the interior.
\end{proof}

\subsection{Parametric phase, polar angle, and action area}
\label{sec:ii-elipse-fase-area}

\begin{proposition}[Areal law and two angular coordinates]
\label{prop:ii-elipse-area}
The phase \(\theta\) of \(\gamma\) and its continuous polar angle
\(\phi\), lifted from zero, satisfy
\begin{equation}
 \phi(\theta)=\arg(a_S\cos\theta+i b_S\sin\theta),\qquad
 \frac{d\phi}{d\theta}
 =\frac{a_Sb_S}{a_S^2\cos^2\theta+b_S^2\sin^2\theta}>0.
 \label{eq:ii-elipse-fases}
\end{equation}
In a chart where both tangents are defined,
\(\tan\phi=(b_S/a_S)\tan\theta\). The oriented swept area is
\begin{equation}
 \mathcal A(\theta)=\frac12\int_0^\theta(QP'-PQ')\,dt
 =\hbar\theta,\qquad
 \mathcal A(1)=\hbar,\qquad \mathcal A(2\pi)=h.
 \label{eq:ii-elipse-area-fase}
\end{equation}
\end{proposition}

\begin{proof}
The derivative of the argument of a nonzero vector is
\((QP'-PQ')/(Q^2+P^2)\). Here the numerator is
\(a_Sb_S=2\hbar\), proving the derivative and monotonicity.
The quotient \(P/Q\) gives the tangent relation; the continuous
lift fixes the quadrants and the number of turns. The areal integral
is \(a_Sb_S\theta/2=\hbar\theta\), and
\(h=2\pi\hbar\) gives the final identity.
\end{proof}

Parametric phase does not generally coincide with polar angle. One radian
of phase sweeps \(\hbar\), but is not thereby identified with a circular
arc or a sector of unit polar angle. In polar coordinates,
the same law is written
\(d\mathcal A=\rho_{\partial}(\phi)^2d\phi/2\).
With the orientation of \(\gamma\) used here,
\(\oint P\,dQ=-h\): positive area is \(h\), whereas the
integral of that one-form preserves its sign. This agrees with the
distinction between symplectic and antisymplectic transports in the
preceding section.

\begin{figure}[htbp]
 \centering
 % Computed geometry, not focal coordinates chosen by eye.
% Illustrative algebraic case xi=1/3; not an evaluation of the HMT constants.
\begingroup
\pgfmathsetmacro{\elXi}{1/3}
\pgfmathsetmacro{\elEta}{0.5*ln((1+\elXi)/(1-\elXi))}
\pgfmathsetmacro{\elA}{exp(\elEta/2)}
\pgfmathsetmacro{\elB}{exp(-\elEta/2)}
\pgfmathsetmacro{\elF}{sqrt(\elA*\elA-\elB*\elB)}
\pgfmathsetmacro{\elTheta}{30}
\pgfmathsetmacro{\elX}{\elA*cos(\elTheta)}
\pgfmathsetmacro{\elY}{\elB*sin(\elTheta)}
\pgfmathsetmacro{\elPolar}{atan(\elY/\elX)}
\begin{tikzpicture}[x=2.18cm,y=2.18cm,font=\small,
 every node/.style={inner sep=2pt},>=Stealth]
 \begin{scope}
  \node[font=\bfseries] at (0,1.38) {Shape and orientation};
  \fill[accent!12] (0,0)--(\elA,0)
    plot[domain=0:30,samples=50] ({\elA*cos(\x)},{\elB*sin(\x)})--cycle;
  \draw[ink,thick] (0,0) ellipse[x radius=\elA,y radius=\elB];
  \draw[->,black!55] (-1.42,0)--(1.48,0) node[right]{\(u\)};
  \draw[->,black!55] (0,-1.05)--(0,1.08) node[above]{\(v\)};
  \draw[accent,thick,-{Stealth}] (\elA,0)
    plot[domain=0:30,samples=50] ({\elA*cos(\x)},{\elB*sin(\x)});
  \draw[accent,thick] (0,0)--(\elX,\elY);
  \node[below left,font=\scriptsize] at (0,0) {\(O\)};
  \draw[accent] (.35,0) arc[start angle=0,end angle=\elPolar,radius=.35];
  \node[accent] at (.47,.065) {\(\phi\)};
  \fill[accent] (\elX,\elY) circle[radius=1.4pt];
  \node[above right] at (\elX,\elY) {\(X(\theta)\)};
  \node[accent,align=center,font=\scriptsize] at (1.26,.18) {\(\theta=\pi/6\)};
  \draw[dashed,black!45] (-\elF,0)--(\elX,\elY);
  \draw[dashed,black!45] (\elF,0)--(\elX,\elY);
  \fill[ink] (-\elF,0) circle[radius=1.5pt] node[below]{\(F_-\)};
  \fill[ink] (\elF,0) circle[radius=1.5pt] node[below]{\(F_+\)};
  \draw[<->,black!60] (-\elA,-1.04)--(0,-1.04)
    node[midway,below,font=\scriptsize]{\(a_S/\sqrt{2\hbar}\)};
  \draw[<->,black!60] (-1.32,0)--(-1.32,\elB)
    node[midway,left,font=\scriptsize]{\(\frac{b_S}{\sqrt{2\hbar}}\)};
  \node[align=center,font=\scriptsize] at (0,-1.55)
    {\(\rho_{\partial}(\phi)/\sqrt{2\hbar}=|OX|\)\\
     \(f_S^2=a_S^2-b_S^2\), \quad total area \(h\)};
 \end{scope}
 \begin{scope}[shift={(3.65,0)}]
  \node[font=\bfseries] at (0,1.38) {Reciprocal shape};
  \draw[densely dashed,black!25] (0,0) ellipse[x radius=\elA,y radius=\elB];
  \draw[ink,thick] (0,0) ellipse[x radius=\elB,y radius=\elA];
  \draw[->,black!55] (-1.28,0)--(1.38,0) node[right]{\(u\)};
  \draw[->,black!55] (0,-1.30)--(0,1.26) node[right]{\(v\)};
  \fill[ink] (0,\elF) circle[radius=1.5pt] node[right]{\(F'_+\)};
  \fill[ink] (0,-\elF) circle[radius=1.5pt] node[right]{\(F'_-\)};
  \node[align=center,font=\scriptsize] at (0,-1.55)
    {\(\eta\mapsto-\eta\), \quad \(a_S\leftrightarrow b_S\)\\
     \(R_{\rm el}\mapsto R_{\rm el}^{-1}\), \quad \(a_Sb_S=2\hbar\)};
 \end{scope}
 \draw[->,accent,thick] (1.60,.80)--(2.17,.80)
   node[midway,above,font=\scriptsize]{exchange};
\end{tikzpicture}
\endgroup

 \caption{Semiaxes, foci, and reciprocity of the action ellipse in the
 normalized coordinates \(u=Q/\sqrt{2\hbar}\),
 \(v=P/\sqrt{2\hbar}\). The marked sector has parametric phase
 \(\theta=\pi/6\), different from its polar angle \(\phi\); its action
 area is \(\hbar\pi/6\). The figure uses the algebraic case
 \(\xi=1/3\) to make the relations legible: it does not assign that value
 to the HMT angular output. The foci are computed from the semiaxes.
 On the right, exchanging axes preserves total area \(h\)
 and inverts the dimensionless modular radius.}
 \label{fig:ii-elipse-calculada}
\end{figure}

For the control case of the figure, \(\xi=1/3\), the normalized
coordinates are exactly
\[
 \eta=\tfrac12\log2,\quad
 \frac{a_S}{\sqrt{2\hbar}}=2^{1/4}=1.1892071150\ldots,\quad
 \frac{b_S}{\sqrt{2\hbar}}=
 \frac{f_S}{\sqrt{2\hbar}}=R_{\rm el}=2^{-1/4}=0.8408964153\ldots.
\]
These are regression numbers for the drawing, not values assigned to the angular
generator. The local certificate also checks the circular shape and
the orientations \(\xi=\pm1/3,\pm3/5\).

\subsection{Shape reciprocity and interior distance}

The rectangular representation recovers
\begin{equation}
 R_{\rm el}=\sqrt{b_S/a_S}=e^{-\eta/2},\qquad
 R_{\rm el}^4=\frac{1-\xi}{1+\xi},\qquad
 \xi=\frac{1-R_{\rm el}^4}{1+R_{\rm el}^4}.
 \label{eq:ii-elipse-reciproca}
\end{equation}
The proof consists in substituting
\(2\eta=\log((1+\xi)/(1-\xi))\). Exchanging axes
sends \(\eta\) to \(-\eta\), \(R_{\rm el}\) to
\(R_{\rm el}^{-1}\), and \(\tau=iR_{\rm el}^2\) to \(-1/\tau\).
The product \(a_Sb_S\) is preserved. The radius \(R_{\rm el}\) is dimensionless;
the polar radius \(\rho_{\partial}\) has square-root-of-action units.

The interior also admits the Hilbert projective distance. If
\(z_-,x,y,z_+\) are ordered on a chord of the ellipse, define
\[
 d_H(x,y)=\frac12\log
 \frac{|y-z_-|\,|x-z_+|}{|x-z_-|\,|y-z_+|}.
\]
It is a dimensionless distance: the units cancel in the cross-ratio.
The normalization \((Q,P)\mapsto(Q/a_S,P/b_S)\) takes the ellipse
to the unit disk and preserves that ratio, because lengths on a
single line are multiplied by a common factor. It thus transports the metric
to the Beltrami--Klein model. Neither this interior distance nor the distance
between foci is identified with an atomic length.

\subsection{Bohr radius: a consequence of the action scale}
\label{sec:ii-elipse-bohr}

In the electrodynamic realization, the Bohr scale is expressed as
\(a_{\rm B}=\hbar/(m_ec\alpha)\) \cite{codata2022}. In the chain
of the present article, this reading follows the constructions of
action, vacuum, fine structure, and electronic mass. Here
\(\hbar,m_e,c,\alpha>0\) are received in the same dimensional realization.
\begin{proposition}[Areal expression of the Bohr scale]
\label{prop:ii-elipse-bohr}
In that realization,
\begin{equation}
 a_{\rm B}=\frac{a_Sb_S}{2m_ec\alpha}
 =\frac{\mathcal A(2\pi)}{2\pi m_ec\alpha}
 =\frac{\bar\lambda_C(m_e)}{\alpha},\qquad
 \bar\lambda_C(m_e)=\frac{\hbar}{m_ec}.
 \label{eq:ii-elipse-bohr}
\end{equation}
\end{proposition}
\begin{proof}
The identities \(a_Sb_S=2\hbar\) and
\(\mathcal A(2\pi)=2\pi\hbar\), substituted into the Bohr
expression, give the first two members. The definition of
\(\bar\lambda_C(m_e)\) gives the third.
\end{proof}

The numerator has action dimension, and \(m_ec\) has momentum dimension;
the quotient is a length. This composition uses the product of
the semiaxes and the other three quantities: it does not assert that Bohr is a semiaxis,
a focus, a Hilbert distance, or the modular radius. Nor does it equate
electronic mass with the circular quantum of order 108. The notation
\(a_{\rm B}\) avoids confusing this length with the normalization
\(a_0\) of the area operator.

\subsection{Inductive--capacitive realization and constitutive continuity}
\label{sec:ii-elipse-lc}

The electrodynamic bridge of the ellipse preserves an additional chart:
a clock \(\omega>0\) and a reference impedance \(Z_*>0\).
The construction recovered from the corpus defines
\begin{equation}
 L_\eta=\frac{Z_*}{\omega}e^{-\eta},\qquad
 C_\eta=\frac{1}{Z_*\omega}e^\eta,
 \qquad Z_\eta=\sqrt{L_\eta/C_\eta}.
 \label{eq:ii-elipse-lc}
\end{equation}
Here \(C_\eta\) is capacitance, distinct from the counter-angle \(C^*\).
\begin{proposition}[Dynamical product and constitutive quotient]
\label{prop:ii-elipse-lc}
In the preceding chart,
\begin{equation}
 L_\eta C_\eta=\omega^{-2},\qquad
 \frac{Z_\eta}{Z_*}=e^{-\eta}
 =\frac{b_S}{a_S}=R_{\rm el}^{2}.
 \label{eq:ii-elipse-impedancia}
\end{equation}
With \(Q=\sqrt{Z_*}\,q\), \(P=\Phi/\sqrt{Z_*}\), the energy
\(H_{LC}=q^2/(2C_\eta)+\Phi^2/(2L_\eta)\) takes the form
\[
 H_{LC}=\frac\omega2(e^{-\eta}Q^2+e^\eta P^2).
\]
On \(\gamma(\theta)\), its electric and magnetic components are
\(U_E=\hbar\omega\cos^2\theta\) and
\(U_B=\hbar\omega\sin^2\theta\); their sum is \(\hbar\omega\).
\end{proposition}
\begin{proof}
The product in \eqref{eq:ii-elipse-lc} cancels the exponents and
\(Z_*\); the quotient is \(Z_*^2e^{-2\eta}\). Taking the positive
root and using \eqref{eq:ii-elipse-reciproca} gives
\eqref{eq:ii-elipse-impedancia}. Substituting \(q=Q/\sqrt{Z_*}\)
and \(\Phi=\sqrt{Z_*}P\) into \(H_{LC}\) gives its balanced form.
Finally, the semiaxes of \(\gamma\) cancel the factors
\(e^{\mp\eta}\), and \(\cos^2\theta+\sin^2\theta=1\).
\end{proof}

This bridge distinguishes shape, clock, and transducer. The choice
of dynamical direction must agree with the symplectic convention: for
the equations \(\dot Q=\partial_PH\),
\(\dot P=-\partial_QH\), the preceding parametrization satisfies
\(\dot\theta=-\omega\). The energy identities do not depend
on reversing that path.

The vacuum resolvent reader acts on the already produced channels
\(q_\pm=e^{-(A_{\rm rad}\pm C^*_{\rm rad})}\) and incidence
\(90/120\). Its responses
\(r_\pm=(1-q_\pm^{90})/(1-q_\pm^{120})\) express
\(\widehat Z=r_-/r_+\) and \(\widehat c=(r_+r_-)^{-1}\).
These are spectral responses, not geometric radii. Identity
\eqref{eq:ii-elipse-impedancia} fixes the LC realization here;
identifying it with the constitutive vacuum impedance preserves the map
between the two readers and their charts. Their values are not equated merely because
they share the angular antecedent.
